\section{Related Work}
\label{sec:related}

\paragraph{Token reduction in ViTs.}
DynamicViT~\cite{rao2021dynamicvit} learns a lightweight prediction module that drops tokens at fixed depths and requires joint training with a distillation loss. A-ViT~\cite{yin2022avit} halts individual tokens adaptively. EViT~\cite{liang2022evit} ranks tokens by the class-token attention and fuses the inattentive remainder into a single token between attention and MLP. ToMe~\cite{bolya2023tome} is training-free: it bipartite-matches tokens by key similarity and merges the most similar pairs at every layer, tracking token mass through proportional attention. Zero-TPrune~\cite{wang2024zerotprune} and Token Cropr~\cite{bergner2025cropr} extend training-free pruning to attention-graph importance and to dense prediction tasks respectively. All of these are developed and evaluated on ImageNet-scale data. Our adapters are controlled simplifications of EViT and ToMe applied after a complete block (\cref{sec:method}); we do not reproduce their exact schedules, and we compare them to the two obvious baselines, random and attention-ranked deletion, under identical budgets.

\paragraph{Efficient inference for medical imaging.}
Zhou \etal~\cite{zhou2023tokensparse} sparsify tokens for medical segmentation, and Byun \etal~\cite{byun2026adaptive} study adaptive inference for medical ViT classifiers. Both select operating points on validation data; to our knowledge no prior work on medical token reduction locks its configurations before test evaluation or reports paired non-inferiority bounds~\cite{walker2011noninferiority}, which we argue is necessary because validation and test sets in MedMNIST are small and, for \pathmnist, come from different clinical centres~\cite{yang2023medmnist}.

\paragraph{Parameter-efficient fine-tuning.}
Raghu \etal~\cite{raghu2019transfusion} showed that ImageNet transfer to medical images often yields most of its benefit from the early layers, so training only the head or the last blocks can be competitive. LoRA~\cite{hu2022lora} adds low-rank updates to attention projections and is now standard for language models; its behaviour on small medical vision datasets, and in particular as a recovery mechanism after token pruning, is less studied. We compare full, partial (last-$k$-block), linear-probe and LoRA adaptation under a common efficiency accounting.
